% =====================================================================
% cap1.tex -- Capitulo 1: A Instalacao (The Facility)
% Estrutura de secoes 1.1 a 1.8 conforme NUREG-1537, Parte 1
% (Format and Content) \citep{nureg1537p1} e criterios de aceitacao
% da NUREG-1537, Parte 2 (Standard Review Plan) \citep{nureg1537p2}.
% =====================================================================
\chapter{A Instala\c{c}\~ao}
\nusecbreaknopage{1.1}{Se\c{c}\~ao 1.1}
\chapterbanner{1}{A Instala\c{c}\~ao}

Este cap\'itulo apresenta uma introdu\c{c}\~ao ao Relat\'orio Preliminar de
An\'alise de Seguran\c{c}a (RPAS) e \`a instala\c{c}\~ao, seguindo o formato
sugerido pela NUREG-1537, Parte~1, \emph{``Guidelines for Preparing and
Reviewing Applications for the Licensing of Non-Power Reactors: Format and
Content''} \citep{nureg1537p1}, e revisado em rela\c{c}\~ao aos crit\'erios de
aceita\c{c}\~ao da NUREG-1537, Parte~2, \emph{``Guidelines for Preparing and
Reviewing Applications for the Licensing of Non-Power Reactors: Standard
Review Plan and Acceptance Criteria''} \citep{nureg1537p2}. Este cap\'itulo \'e uma vis\~ao geral, ou resumo
executivo, dos temas tratados em detalhe nos cap\'itulos subsequentes do
RPAS; ele n\~ao substitui a discuss\~ao completa e as an\'alises
apresentadas naqueles cap\'itulos.
%\pagebreak

\begin{notaregulatoria}{Sobre esta proposta}
Este documento \'e um \textbf{modelo/proposta} (template) para o
Cap\'itulo~1 do RPAS, com trechos entre colchetes \texttt{[...]} a
serem preenchidos com os dados reais da instala\c{c}\~ao. A estrutura de
se\c{c}\~oes (\ref{sec:1.1} a \ref{sec:1.8}) segue integralmente a NUREG-1537, Parte~1
\citep{nureg1537p1}, que \'e voltada a reatores n\~ao-de-pot\^encia
(reatores de pesquisa e multiprop\'osito).
\end{notaregulatoria}

% =====================================================================
\section{Introdu\c{c}\~ao}
\label{sec:1.1}

[Nome do Requerente], uma [universidade / \'org\~ao governamental /
instituto de pesquisa / empresa], apresenta este Cap\'itulo~1 do
Relat\'orio Preliminar de An\'alise de Seguran\c{c}a (RPAS) em apoio a
uma solicita\c{c}\~ao de [permiss\~ao de constru\c{c}\~ao e licen\c{c}a de
opera\c{c}\~ao] para a [Nome da
Instala\c{c}\~ao/Reator], um [tipo de reator, p.ex., reator a \'agua
pressurizada sem boro sol\'uvel (SBF-PWR) / reator de pesquisa
multiprop\'osito tipo piscina aberta] com pot\^encia t\'ermica nominal de
[valor] MW(t) [e, se aplic\'avel, capacidade de pulso at\'e [valor]
MW(t) de pico], a ser instalado em [local/cidade, estado, pa\'is]. O uso
pretendido da instala\c{c}\~ao \'e [pesquisa, produ\c{c}\~ao de
rad\i s\'otopos, irradia\c{c}\~ao de materiais, pesquisa com feixes de
n\^eutrons, ensino e treinamento e/ou gera\c{c}\~ao de eletricidade,
conforme aplic\'avel].

As principais caracter\'isticas de seguran\c{c}a inerentes ou passivas do
projeto incluem [p.ex., coeficiente de temperatura do moderador
fortemente negativo, remo\c{c}\~ao do calor de decaimento por
circula\c{c}\~ao natural, sistemas passivos de resfriamento do n\'ucleo e
controle de reatividade sem boro sol\'uvel, apoiado em venenos
queim\'aveis e barras de controle]. Caracter\'isticas de projeto \'unicas
que distinguem esta instala\c{c}\~ao de reatores n\~ao-de-pot\^encia ou de
pot\^encia previamente licenciados incluem [listar, p.ex., elimina\c{c}\~ao
da fun\c{c}\~ao de bora\c{c}\~ao do sistema de controle qu\'imico e de
volume, n\'ucleo compacto com retroalimenta\c{c}\~ao de reatividade
negativa refor\c{c}ada, ou miss\~ao combinada de pesquisa/irradia\c{c}\~ao].
Esses temas s\~ao tratados integralmente nos Cap\'itulos~2 a~17 deste
RPAS, aos quais esta se\c{c}\~ao remete ao longo do texto.

\subsection*{Enquadramento regulat\'orio do RMB}
\label{sec:1.1.enquadramento}

Antes de entrar no detalhamento dos itens ~\ref{sec:1.1} a ~\ref{sec:1.8}, cabe situar onde o
RMB se encaixa nas categorias que a NUREG-1537 \citep{nureg1537p1,nureg1537p2} e o \emph{Interim Staff
Guidance Augmenting NUREG-1537}
\citep[NPR-ISG-2011-002;][]{isgnureg1537p1,isgnureg1537p2} -- aqui referido, por
brevidade, como \textbf{ISG-NUREG-1537} -- que a complementa efetivamente
reconhecem -- porque nenhuma das duas normas usa a express\~ao ``reator
multiprop\'osito'' como categoria regulat\'oria pr\'opria.

A NUREG-1537 trabalha com uma divis\~ao bin\'aria dentro do guarda-chuva
de reator n\~ao-de-pot\^encia: de um lado o reator de pesquisa
(\emph{research reactor}), de outro a instala\c{c}\~ao de teste
(\emph{testing facility}), esta \'ultima sujeita a exig\^encias
adicionais -- declara\c{c}\~ao de impacto ambiental, audi\^encia de
licenciamento e revis\~ao pelo \emph{Advisory Committee on Reactor
Safeguards} (ACRS). O crit\'erio que separa as duas categorias \'e
objetivo, n\~ao nominal: o 10~CFR~50.2 -- \citep{cfr10p50} -- enquadra como \emph{testing
facility} qualquer reator com pot\^encia t\'ermica acima de 1~MW(t) e com
instala\c{c}\~ao experimental \emph{in-core} de se\c{c}\~ao transversal
superior a 16~polegadas quadradas. Com pot\^encia de projeto de 30~MW(t)
e m\'ultiplos canais de irradia\c{c}\~ao \emph{in-core} e
\emph{in-reflector}, o RMB tende a se enquadrar nessa categoria mais
rigorosa -- n\~ao por ser ``multiprop\'osito'', mas por pot\^encia e por
ter instala\c{c}\~oes experimentais internas ao n\'ucleo acima do limiar
do regulamento.
%esse enquadramento deve ser confirmado formalmente com o
%\'org\~ao regulador nacional antes %da submiss\~ao deste RPAS.

O ISG-NUREG-1537 acrescenta uma segunda camada a
essa classifica\c{c}\~ao. Ele foi escrito para cobrir, entre outros
cen\'arios, exatamente o do RMB: um reator \textbf{heterog\^eneo}
(n\'ucleo com elementos combust\'iveis tipo placa, n\~ao solu\c{c}\~ao
homog\^enea) que irradia alvos de ur\^anio pouco enriquecido (LEU) para
produ\c{c}\~ao de molibd\^enio-99, com o material irradiado
posteriormente processado em uma instala\c{c}\~ao de separa\c{c}\~ao
radioqu\'imica associada. Para esse cen\'ario -- que o pr\'oprio ISG-NUREG-1537
descreve como produ\c{c}\~ao de Mo-99 por fiss\~ao de material nuclear
especial em alvos LEU irradiados em reatores n\~ao-de-pot\^encia
existentes ou novos --, o tratamento \'e expl\'icito: o reator em si
continua coberto pelo escopo padr\~ao da NUREG-1537, com orienta\c{c}\~ao
adicional m\'inima; j\'a a instala\c{c}\~ao onde ocorre a separa\c{c}\~ao
do radiois\'otopo passa a ser tratada, para fins de licenciamento, como
uma instala\c{c}\~ao de produ\c{c}\~ao (\emph{production facility}) sob o
10~CFR Parte~50 (\citealp{cfr10p50}), sujeita aos cap\'itulos e crit\'erios de aceita\c{c}\~ao
que o ISG-NUREG-1537 acrescenta especificamente para esse tipo de instala\c{c}\~ao.

Na pr\'atica, isso separa o RMB em dois objetos de licenciamento com
tratamento regulat\'orio distinto, ainda que fisicamente cont\'iguos: o
pr\'edio do reator, onde a NUREG-1537 domina quase sem altera\c{c}\~ao
(com os complementos ``a1'' do ISG-NUREG-1537 para reator heterog\^eneo, quando
aplic\'aveis aos Cap\'itulos~4, 5, 6, 8 e~9), e a instala\c{c}\~ao de
produ\c{c}\~ao de radiois\'otopos, onde o ISG-NUREG-1537 \'e a refer\^encia
prim\'aria. Essa distin\c{c}\~ao \'e consistente com a estrat\'egia de
licenciamento faseado j\'a registrada no anexo deste cap\'itulo (ver
``Anexo -- Estrat\'egia de Licenciamento Faseado''): o escopo inicial
deste RPAS, limitado ao pr\'edio do reator, corresponde ao objeto onde a
NUREG-1537 base j\'a basta; a instala\c{c}\~ao de produ\c{c}\~ao de Mo-99,
quando entrar em escopo, dever\'a ser tratada com o apoio direto do
ISG-NUREG-1537, mapeado no \apref{apA}{A} deste cap\'itulo.

\begin{tabelaaceitacao}{\ref{sec:1.1}}{Introdu\c{c}\~ao}
1 & O prop\'osito do RPAS deve estar claramente declarado. \\
2 & O requerente deve estar identificado (nome, natureza jur\'idica). \\
3 & A localiza\c{c}\~ao, o prop\'osito e o uso da instala\c{c}\~ao devem
     estar brevemente descritos. \\
4 & As caracter\'isticas b\'asicas da instala\c{c}\~ao que afetam as
     considera\c{c}\~oes de licenciamento devem estar brevemente
     discutidas. \\
5 & As caracter\'isticas de projeto ou de localiza\c{c}\~ao voltadas a
     endere\c{c}ar preocupa\c{c}\~oes b\'asicas de seguran\c{c}a devem
     estar delineadas. \\
6 & Caracter\'isticas de seguran\c{c}a \'unicas do projeto, diferentes de
     instala\c{c}\~oes n\~ao-de-pot\^encia previamente licenciadas, devem
     estar destacadas. \\
\end{tabelaaceitacao}

\begin{tabelaresumo}{\ref{sec:1.1}}{Introdu\c{c}\~ao}
Identifica\c{c}\~ao do requerente & Nome, natureza jur\'idica (universidade,
  \'org\~ao governamental, instituto, empresa) e, quando aplic\'avel, o
  arranjo de propriedade/opera\c{c}\~ao da instala\c{c}\~ao. \\
Prop\'osito e uso pretendido & Miss\~ao declarada do reator (pesquisa,
  irradia\c{c}\~ao, produ\c{c}\~ao de is\'otopos, ensino, gera\c{c}\~ao de
  pot\^encia). \\
Localiza\c{c}\~ao geogr\'afica & Local, munic\'ipio/estado/pa\'is do
  s\'itio proposto. \\
Tipo e pot\^encia do reator & Tecnologia (p.ex. PWR sem boro sol\'uvel,
  piscina aberta), pot\^encia t\'ermica nominal e, se aplic\'avel,
  capacidade de pulso. \\
Caracter\'isticas de seguran\c{c}a inerentes/passivas & Elementos de
  seguran\c{c}a intr\'insecos ao projeto que reduzem a depend\^encia de
  a\c{c}\~ao ativa ou humana. \\
Caracter\'isticas \'unicas de projeto & Aspectos que diferenciam esta
  instala\c{c}\~ao de projetos j\'a licenciados e que merecer\~ao aten\c{c}\~ao
  espec\'ifica do revisor. \\
\end{tabelaresumo}

\begin{tabelaapendices}{\ref{sec:1.1}}{Introdu\c{c}\~ao}
Cap\'itulos 2--17 do RPAS & Confirmar que cada tema mencionado nesta
  se\c{c}\~ao (local, tipo de reator, pot\^encia, seguran\c{c}a passiva)
  \'e de fato detalhado e \'e consistente com o cap\'itulo correspondente. &
  Pendente \\
NUREG-1537, Parte~1, \S1.1 \citep{nureg1537p1} & Verificar aderência ao
  conte\'udo m\'inimo exigido para a introdu\c{c}\~ao. & Pendente \\
NUREG-1537, Parte~2, \S1.1 \citep{nureg1537p2} & Verificar os crit\'erios
  de aceita\c{c}\~ao aplic\'aveis antes da submiss\~ao. & Pendente \\
\end{tabelaapendices}

\begin{avaliacaointerim}
O ISG-NUREG-1537 amplia o escopo desta se\c{c}\~ao:
onde a NUREG-1537 fala em caracter\'isticas de seguran\c{c}a inerentes ou
passivas do reator, o ISG-NUREG-1537 passa a exigir que as caracter\'isticas
inerentes/passivas da instala\c{c}\~ao de produ\c{c}\~ao de radiois\'otopos
associada tamb\'em sejam descritas aqui, quando aplic\'avel ao escopo do
RPAS. N\~ao h\'a altera\c{c}\~ao na estrutura da se\c{c}\~ao nem novo
crit\'erio de aceita\c{c}\~ao formal -- apenas amplia\c{c}\~ao do escopo
textual, j\'a refletida no \apref{apA}{A}. Como o escopo inicial deste RPAS
est\'a limitado ao pr\'edio do reator (ver Anexo -- Estrat\'egia de
Licenciamento Faseado), essa amplia\c{c}\~ao passa a valer integralmente
apenas quando a instala\c{c}\~ao de produ\c{c}\~ao de Mo-99 for
incorporada em fase posterior.
\end{avaliacaointerim}

% =====================================================================
\nusecbreak{1.2}{Se\c{c}\~ao 1.2}
\section{Resumo e Conclus\~oes sobre as Principais Considera\c{c}\~oes de
Seguran\c{c}a}
\label{sec:1.2}

Esta se\c{c}\~ao apresenta os crit\'erios de seguran\c{c}a, as principais
considera\c{c}\~oes de seguran\c{c}a e as conclus\~oes decorrentes para a
instala\c{c}\~ao, incluindo discuss\~oes breves sobre: as consequ\^encias
da opera\c{c}\~ao e uso do reator e os m\'etodos empregados para garantir
a seguran\c{c}a; as considera\c{c}\~oes de seguran\c{c}a que influenciaram
a escolha do s\'itio, do tipo de reator e combust\'ivel, da pot\^encia
t\'ermica e do tipo de edifica\c{c}\~ao; as caracter\'isticas de
seguran\c{c}a inerentes ou passivas projetadas para proteger a sa\'ude e
seguran\c{c}a do p\'ublico, da equipe e do meio ambiente; as bases de
projeto de sistemas e componentes que promovem a opera\c{c}\~ao e o
desligamento seguros; e os acidentes potenciais na instala\c{c}\~ao,
incluindo o acidente hipot\'etico m\'aximo (MHA) [e, quando aplic\'avel,
outros eventos-base de projeto limitantes], bem como as caracter\'isticas
de projeto que os previnem ou mitigam.

N\~ao se espera que nenhuma exposi\c{c}\~ao decorrente da opera\c{c}\~ao
normal exceda os limites do [10~CFR Parte~20 \citep{cfr10p20} / regulamento
nacional de radioprote\c{c}\~ao aplic\'avel], e a instala\c{c}\~ao \'e
projetada e ser\'a operada de acordo com o princ\'ipio ALARA (t\~ao baixo
quanto razoavelmente exequ\'ivel). Essas discuss\~oes constituem uma
vis\~ao geral; as an\'alises detalhadas encontram-se nos
Cap\'itulos~3 a~13 deste RPAS.
\begin{tabelaaceitacao}{\ref{sec:1.2}}{Resumo das considera\c{c}\~oes de seguran\c{c}a}
1 & Caracter\'isticas de projeto suficientes devem estar inclu\'idas para
     proteger a sa\'ude e a seguran\c{c}a do p\'ublico. \\
2 & Nenhuma exposi\c{c}\~ao decorrente de opera\c{c}\~ao normal deve
     exceder os limites regulat\'orios aplic\'aveis, observado o
     princ\'ipio ALARA. \\
3 & Os acidentes potenciais devem estar brevemente discutidos. \\
4 & Todos os modos de opera\c{c}\~ao e eventos capazes de causar liberações
     radiol\'ogicas significativas e exposi\c{c}\~ao do p\'ublico devem
     estar discutidos. \\
\end{tabelaaceitacao}

\begin{tabelaresumo}{\ref{sec:1.2}}{Resumo das considera\c{c}\~oes de seguran\c{c}a}
Crit\'erios de seguran\c{c}a declarados & Base regulat\'oria e princ\'ipios
  de projeto adotados (defesa em profundidade, ALARA, margens de
  seguran\c{c}a). \\
Considera\c{c}\~oes que influenciaram o projeto & Justificativas de
  escolha do s\'itio, tipo de reator/combust\'ivel, pot\^encia e
  edifica\c{c}\~ao. \\
Caracter\'isticas inerentes/passivas & Mecanismos que independem de a\c{c}\~ao
  ativa (retroalimenta\c{c}\~ao de reatividade, remo\c{c}\~ao passiva de
  calor). \\
Acidente hipot\'etico m\'aximo (MHA) & Cen\'ario acidental mais restritivo
  considerado no dimensionamento das barreiras de seguran\c{c}a. \\
Conclus\~ao de conformidade & Declara\c{c}\~ao de que as doses previstas
  respeitam os limites regulat\'orios em todas as condi\c{c}\~oes
  avaliadas. \\
\end{tabelaresumo}

\begin{tabelaapendices}{\ref{sec:1.2}}{Resumo das considera\c{c}\~oes de seguran\c{c}a}
Cap\'itulo 13 -- An\'alise de Acidentes & Confirmar que o acidente
  hipot\'etico m\'aximo (MHA) citado aqui corresponde \`a an\'alise
  detalhada do Cap\'itulo 13. & Pendente \\
Cap\'itulo 6 -- Caracter\'isticas de Seguran\c{c}a Projetadas & Verificar
  consist\^encia entre as caracter\'isticas passivas resumidas aqui e a
  descri\c{c}\~ao detalhada. & Pendente \\
10~CFR Parte~20 (ou equivalente nacional) \citep{cfr10p20} & Confirmar que
  os limites de dose citados est\~ao corretos e atualizados. & Pendente \\
\end{tabelaapendices}

\begin{avaliacaointerim}
Esta \'e a se\c{c}\~ao em que o ISG-NUREG-1537 introduz a mudan\c{c}a mais
substantiva para o Cap\'itulo~1: al\'em de reformular itens de texto
corrido da NUREG-1537 para incluir explicitamente a instala\c{c}\~ao de
produ\c{c}\~ao de radiois\'otopos, o ISG-NUREG-1537 acrescenta, na Parte~2
(crit\'erios de aceita\c{c}\~ao), uma \'area de revis\~ao inteiramente
nova -- a vis\~ao geral de processo (\emph{process overview}) da
instala\c{c}\~ao de produ\c{c}\~ao, cobrindo a rela\c{c}\~ao entre os
principais processos qu\'imicos/mec\^anicos e o material licenci\'avel
envolvido, a localiza\c{c}\~ao das edifica\c{c}\~oes dos principais
componentes de processo (apoiada em desenhos), e a marca\c{c}\~ao de
eventuais informa\c{c}\~oes propriet\'arias ou sens\'iveis. Diferentemente
das demais se\c{c}\~oes, este n\~ao \'e apenas um ajuste de terminologia:
se e quando a instala\c{c}\~ao de produ\c{c}\~ao de Mo-99 entrar no
escopo deste RPAS, esta se\c{c}\~ao vai exigir um par\'agrafo pr\'oprio de
vis\~ao geral de processo, com n\'ivel de detalhe consistente com o
resumo da an\'alise de seguran\c{c}a integrada (ISA) e com o
Cap\'itulo~13. Essa necessidade j\'a est\'a sinalizada no \apref{apA}{A}.
\end{avaliacaointerim}

% =====================================================================
\nusecbreak{1.3}{Se\c{c}\~ao 1.3}
\section{Descri\c{c}\~ao Geral}
\label{sec:1.3}

Esta se\c{c}\~ao descreve brevemente a instala\c{c}\~ao do reator,
incluindo: (1) a localiza\c{c}\~ao geogr\'afica; (2) as principais
caracter\'isticas do s\'itio; (3) os principais crit\'erios de projeto,
caracter\'isticas operacionais e sistemas de seguran\c{c}a; (4) quaisquer
caracter\'isticas de seguran\c{c}a projetadas (\textit{engineered safety
features}); (5) sistemas de instrumenta\c{c}\~ao, controle e el\'etricos;
(6) o sistema de refrigera\c{c}\~ao do reator e demais sistemas
auxiliares; (7) as provis\~oes de gerenciamento de rejeitos radioativos e
prote\c{c}\~ao radiol\'ogica; e (8) as instala\c{c}\~oes e capacidades
experimentais [ou, no caso de um projeto de reator de pot\^encia, os
sistemas do \textit{balance-of-plant}]. O arranjo geral das principais
estruturas e equipamentos \'e indicado por meio de plantas e cortes nas
Figuras~1-[n] a~1-[n]. As caracter\'isticas de seguran\c{c}a de maior
interesse -- tais como [caracter\'isticas incomuns do s\'itio, o projeto
de cont\^enc\~ao/confinamento, aspectos in\'editos do projeto do reator ou
instala\c{c}\~oes experimentais \'unicas] -- s\~ao brevemente destacadas a
seguir; a discuss\~ao e a an\'alise completas encontram-se nos
cap\'itulos subsequentes deste RPAS, aos quais esta se\c{c}\~ao remete.
\begin{tabelaaceitacao}{\ref{sec:1.3}}{Descri\c{c}\~ao geral da instala\c{c}\~ao}
1 & Devem estar brevemente descritos: localiza\c{c}\~ao geogr\'afica,
     caracter\'isticas do s\'itio, crit\'erios de projeto, sistemas de
     seguran\c{c}a, instrumenta\c{c}\~ao/controle/el\'etrica, sistemas de
     refrigera\c{c}\~ao e auxiliares, gerenciamento de rejeitos e
     prote\c{c}\~ao radiol\'ogica, e instala\c{c}\~oes experimentais. \\
2 & O arranjo geral das principais estruturas e equipamentos deve estar
     indicado por plantas e cortes (desenhos de eleva\c{c}\~ao). \\
3 & As caracter\'isticas de seguran\c{c}a de maior interesse devem estar
     brevemente identificadas. \\
4 & Caracter\'isticas incomuns do s\'itio, da cont\^enc\~ao, do projeto do
     reator ou de instala\c{c}\~oes experimentais \'unicas devem estar
     destacadas. \\
\end{tabelaaceitacao}

\begin{tabelaresumo}{\ref{sec:1.3}}{Descri\c{c}\~ao geral da instala\c{c}\~ao}
S\'itio & Localiza\c{c}\~ao e principais caracter\'isticas f\'isicas,
  demogr\'aficas e ambientais (detalhado no Cap\'itulo 2). \\
Sistemas de seguran\c{c}a e projeto & Vis\~ao geral dos crit\'erios de
  projeto e sistemas de seguran\c{c}a (detalhado nos Cap\'itulos 3, 4 e 6). \\
Instrumenta\c{c}\~ao, controle e el\'etrica & Vis\~ao geral (detalhado nos
  Cap\'itulos 7 e 8). \\
Refrigera\c{c}\~ao e auxiliares & Sistemas de remo\c{c}\~ao de calor e
  suporte (detalhado nos Cap\'itulos 5 e 9). \\
Rejeitos e radioprote\c{c}\~ao & Provis\~oes de gerenciamento e prote\c{c}\~ao
  radiol\'ogica (detalhado no Cap\'itulo 11). \\
Instala\c{c}\~oes experimentais & Capacidades experimentais/irradia\c{c}\~ao
  (detalhado no Cap\'itulo 10), quando aplic\'avel. \\
Desenhos de arranjo geral & Plantas e cortes indicando a disposi\c{c}\~ao
  das principais estruturas e equipamentos. \\
\end{tabelaresumo}

\begin{tabelaapendices}{\ref{sec:1.3}}{Descri\c{c}\~ao geral da instala\c{c}\~ao}
Cap\'itulos 2 a 11 do RPAS & Confirmar consist\^encia entre o resumo desta
  se\c{c}\~ao e o conte\'udo detalhado de cada cap\'itulo referenciado. &
  Pendente \\
Figuras 1-[n] a 1-[n] (plantas e cortes) & Verificar exist\^encia,
  numera\c{c}\~ao e atualiza\c{c}\~ao dos desenhos de arranjo geral. &
  Pendente \\
\end{tabelaapendices}

\begin{avaliacaointerim}
O ISG-NUREG-1537 segue o mesmo padr\~ao de amplia\c{c}\~ao j\'a visto na
Se\c{c}\~ao~\ref{sec:1.1}: a descri\c{c}\~ao geral da instala\c{c}\~ao, tal como
exigida pela NUREG-1537, deve ser expandida para tamb\'em cobrir a
instala\c{c}\~ao de produ\c{c}\~ao de radiois\'otopos, quando esta
estiver no escopo do RPAS. N\~ao h\'a crit\'erio de aceita\c{c}\~ao
adicional espec\'ifico para esta se\c{c}\~ao na Parte~2 do ISG-NUREG-1537, nem
altera\c{c}\~ao estrutural -- apenas a mesma substitui\c{c}\~ao de leitura
(``reator'' passa a significar ``reator e/ou instala\c{c}\~ao de
produ\c{c}\~ao de radiois\'otopos'') j\'a registrada na Se\c{c}\~ao~\ref{sec:1.1}.
\end{avaliacaointerim}

% =====================================================================
\nusecbreak{1.4}{Se\c{c}\~ao 1.4}
\section{Instala\c{c}\~oes e Equipamentos Compartilhados}
\label{sec:1.4}

Esta se\c{c}\~ao descreve brevemente: sistemas e equipamentos
compartilhados com instala\c{c}\~oes n\~ao abrangidas por este RPAS ou
pela licen\c{c}a/certifica\c{c}\~ao (p.ex., sistemas de purifica\c{c}\~ao
de \'agua; suprimento el\'etrico; sistemas de aquecimento, ventila\c{c}\~ao
e ar-condicionado; e o edif\'icio que abriga a sala do reator); qualquer
outro reator, conjunto subcr\'itico, instala\c{c}\~ao de irradia\c{c}\~ao
ou c\'elula quente localizados dentro da estrutura de
confinamento/cont\^enc\~ao ou da \'area restrita a que este RPAS se
aplica; e quaisquer barreiras de seguran\c{c}a ou provis\~oes especiais de
isolamento para as instala\c{c}\~oes e equipamentos compartilhados. [Caso
n\~ao haja instala\c{c}\~oes ou equipamentos compartilhados, isso deve ser
declarado explicitamente.] As descri\c{c}\~oes completas e quaisquer
implica\c{c}\~oes de seguran\c{c}a decorrentes do compartilhamento s\~ao
avaliadas nos cap\'itulos pertinentes deste RPAS, aos quais esta se\c{c}\~ao
remete.
\begin{tabelaaceitacao}{\ref{sec:1.4}}{Instala\c{c}\~oes e equipamentos compartilhados}
1 & A instala\c{c}\~ao deve ser projetada para acomodar todos os usos ou
     falhas das instala\c{c}\~oes compartilhadas sem degrada\c{c}\~ao das
     caracter\'isticas de seguran\c{c}a do reator. \\
2 & O projeto deve evitar condi\c{c}\~oes em que a contamina\c{c}\~ao
     possa se propagar para as instala\c{c}\~oes ou equipamentos
     compartilhados. \\
3 & Onde necess\'ario, as barreiras devem estar brevemente descritas para
     garantir o atendimento aos dois crit\'erios anteriores. \\
\end{tabelaaceitacao}

\begin{tabelaresumo}{\ref{sec:1.4}}{Instala\c{c}\~oes e equipamentos compartilhados}
Sistemas compartilhados & Lista de sistemas/equipamentos compartilhados
  (\'agua, energia el\'etrica, HVAC, edifica\c{c}\~ao) e seus outros
  usu\'arios. \\
Instala\c{c}\~oes co-localizadas & Outros reatores, conjuntos subcr\'iticos
  ou instala\c{c}\~oes de irradia\c{c}\~ao no mesmo pr\'edio/\'area
  restrita. \\
Barreiras e isolamento & Provis\~oes que impedem a propaga\c{c}\~ao de
  contamina\c{c}\~ao ou a perda de fun\c{c}\~ao de seguran\c{c}a por conta
  do compartilhamento. \\
\end{tabelaresumo}

\begin{tabelaapendices}{\ref{sec:1.4}}{Instala\c{c}\~oes e equipamentos compartilhados}
Cap\'itulo 9 -- Sistemas Auxiliares & Confirmar descri\c{c}\~ao detalhada
  dos sistemas compartilhados citados aqui. & Pendente \\
Cap\'itulo 13 -- An\'alise de Acidentes & Verificar se falhas de sistemas
  compartilhados foram avaliadas quanto \`a libera\c{c}\~ao de material
  radioativo. & Pendente \\
\end{tabelaapendices}

\begin{avaliacaointerim}
O ISG-NUREG-1537 trata as Se\c{c}\~oes~\ref{sec:1.4} a~\ref{sec:1.8} em bloco, sem acr\'escimo
espec\'ifico a esta se\c{c}\~ao al\'em da mesma amplia\c{c}\~ao
terminol\'ogica j\'a aplicada \`as se\c{c}\~oes anteriores: onde a
NUREG-1537 menciona ``reator'', o ISG-NUREG-1537 entende que a men\c{c}\~ao
cobre tamb\'em a instala\c{c}\~ao de produ\c{c}\~ao de radiois\'otopos,
quando aplic\'avel. Para esta se\c{c}\~ao em particular, isso significa
que sistemas e equipamentos eventualmente compartilhados entre o pr\'edio
do reator e uma futura instala\c{c}\~ao de produ\c{c}\~ao de Mo-99 devem,
quando essa instala\c{c}\~ao entrar em escopo, ser descritos aqui nos
mesmos termos j\'a usados para outras instala\c{c}\~oes compartilhadas.
N\~ao h\'a mudan\c{c}a estrutural decorrente do ISG-NUREG-1537 nesta se\c{c}\~ao.
\end{avaliacaointerim}

% =====================================================================
\nusecbreak{1.5}{Se\c{c}\~ao 1.5}
\section{Compara\c{c}\~ao com Instala\c{c}\~oes Similares}
\label{sec:1.5}

Esta se\c{c}\~ao descreve brevemente as principais semelhan\c{c}as entre
esta instala\c{c}\~ao e outras instala\c{c}\~oes, particularmente aquelas
licenciadas pela [autoridade regulat\'oria nacional] ou projetadas e
operadas por organiza\c{c}\~oes de pesquisa reconhecidas. S\~ao
comparados os principais par\^ametros de projeto, os sistemas de
seguran\c{c}a do reator, as caracter\'isticas de seguran\c{c}a projetadas
e os sistemas de instrumenta\c{c}\~ao e controle. O hist\'orico
operacional das instala\c{c}\~oes de refer\^encia \'e brevemente citado
para demonstrar a seguran\c{c}a e a confiabilidade da abordagem de
projeto adotada. Caracter\'isticas de projeto, experi\^encia operacional
e testes e experimentos das instala\c{c}\~oes de refer\^encia [p.ex.,
listar instala\c{c}\~oes/tecnologias de refer\^encia] s\~ao utilizados
para embasar as an\'alises apresentadas nos cap\'itulos aplic\'aveis
deste RPAS.
\begin{tabelaaceitacao}{\ref{sec:1.5}}{Compara\c{c}\~ao com instala\c{c}\~oes similares}
1 & As compara\c{c}\~oes devem demonstrar que a instala\c{c}\~ao proposta
     n\~ao excede o envelope de seguran\c{c}a das instala\c{c}\~oes
     similares. \\
2 & Deve haver garantia razo\'avel de que as exposi\c{c}\~oes
     radiol\'ogicas do p\'ublico n\~ao excedem os regulamentos e as
     diretrizes ALARA da instala\c{c}\~ao proposta. \\
\end{tabelaaceitacao}

\begin{tabelaresumo}{\ref{sec:1.5}}{Compara\c{c}\~ao com instala\c{c}\~oes similares}
Instala\c{c}\~oes de refer\^encia & Lista das instala\c{c}\~oes similares
  utilizadas como base comparativa (tipo, pot\^encia, tecnologia,
  status de licenciamento). \\
Par\^ametros comparados & Par\^ametros de projeto, sistemas de seguran\c{c}a,
  instrumenta\c{c}\~ao/controle e desempenho de combust\'ivel. \\
Hist\'orico operacional & Evid\^encias de opera\c{c}\~ao segura e
  confi\'avel nas instala\c{c}\~oes de refer\^encia. \\
Conclus\~ao comparativa & Declara\c{c}\~ao de que o projeto proposto n\~ao
  excede o envelope de seguran\c{c}a j\'a validado. \\
\end{tabelaresumo}

\begin{tabelaapendices}{\ref{sec:1.5}}{Compara\c{c}\~ao com instala\c{c}\~oes similares}
Cap\'itulo 4 -- Descri\c{c}\~ao do Reator & Confirmar base de aceita\c{c}\~ao
  do desempenho do combust\'ivel por compara\c{c}\~ao. & Pendente \\
Cap\'itulo 6 -- Caracter\'isticas de Seguran\c{c}a Projetadas & Confirmar
  base de compara\c{c}\~ao dos sistemas de mitiga\c{c}\~ao de acidentes. &
  Pendente \\
Cap\'itulo 7 -- Instrumenta\c{c}\~ao e Controle & Confirmar base de
  compara\c{c}\~ao de redund\^ancia/diversidade de instrumenta\c{c}\~ao. &
  Pendente \\
\end{tabelaapendices}

\begin{avaliacaointerim}
Diferentemente da amplia\c{c}\~ao gen\'erica aplicada \`as demais
se\c{c}\~oes do bloco \ref{sec:1.4}--\ref{sec:1.8}, a Se\c{c}\~ao~\ref{sec:1.5} recebe do ISG-NUREG-1537 uma
orienta\c{c}\~ao espec\'ifica: a lista de instala\c{c}\~oes de
refer\^encia para compara\c{c}\~ao, hoje limitada a reatores heterog\^eneos
n\~ao-de-pot\^encia, deve incluir tamb\'em instala\c{c}\~oes ou
opera\c{c}\~oes de produ\c{c}\~ao de radiois\'otopos compar\'aveis, quando
aplic\'avel. O ISG-NUREG-1537 lista ainda comparadores espec\'ificos para
reatores homog\^eneos aquosos (AHR) -- LOPO, HYPO, SUPO, TRACY e HRE --
mas esses n\~ao se aplicam ao RMB, que \'e um reator heterog\^eneo tipo
piscina, n\~ao um AHR; a distin\c{c}\~ao fica registrada aqui para evitar
compara\c{c}\~ao indevida. A mudan\c{c}a pr\'atica, quando a
instala\c{c}\~ao de produ\c{c}\~ao de Mo-99 entrar em escopo, \'e incluir
na tabela-resumo desta se\c{c}\~ao ao menos uma refer\^encia de
instala\c{c}\~ao de produ\c{c}\~ao de radiois\'otopos compar\'avel (por
exemplo, produtores de Mo-99 baseados em alvos LEU j\'a licenciados
internacionalmente).
\end{avaliacaointerim}

% =====================================================================
\nusecbreak{1.6}{Se\c{c}\~ao 1.6}
\section{Resumo das Opera\c{c}\~oes}
\label{sec:1.6}

Esta se\c{c}\~ao discute brevemente as opera\c{c}\~oes planejadas do
reator, os programas experimentais e a miss\~ao da instala\c{c}\~ao. As
opera\c{c}\~oes propostas s\~ao importantes para estimar par\^ametros como
tempo total de opera\c{c}\~ao, n\'ivel de pot\^encia, modo de opera\c{c}\~ao
(cont\'inuo, sob demanda ou pulsado) e a quantidade e o tipo de materiais
radioativos subprodutos gerados. [Caso a solicita\c{c}\~ao seja de
renova\c{c}\~ao de licen\c{c}a, esta se\c{c}\~ao deve refletir os planos
operacionais atuais e propostos.] Caso considera\c{c}\~oes de seguran\c{c}a
analisadas em cap\'itulos posteriores deste RPAS limitem o cronograma de
opera\c{c}\~ao do reator, esse fato \'e registrado aqui e referenciado ao
cap\'itulo aplic\'avel e \`as especifica\c{c}\~oes t\'ecnicas
(Cap\'itulo~14).
\begin{tabelaaceitacao}{\ref{sec:1.6}}{Resumo das opera\c{c}\~oes}
1 & O requerente deve demonstrar a consist\^encia das opera\c{c}\~oes
     propostas com as premissas adotadas nos cap\'itulos posteriores do
     RPAS, incluindo os efeitos sobre a integridade do reator e as
     exposi\c{c}\~oes radiol\'ogicas potenciais. \\
2 & O requerente deve demonstrar que a opera\c{c}\~ao proposta do reator
     foi considerada de forma conservadora no projeto e nas an\'alises de
     seguran\c{c}a. \\
3 & (Renova\c{c}\~ao) As opera\c{c}\~oes propostas devem ser consistentes
     com as premissas dos cap\'itulos posteriores do RPAS. \\
\end{tabelaaceitacao}

\begin{tabelaresumo}{\ref{sec:1.6}}{Resumo das opera\c{c}\~oes}
Miss\~ao e programas & Finalidade operacional (pesquisa, irradia\c{c}\~ao,
  produ\c{c}\~ao de is\'otopos, gera\c{c}\~ao) e programas experimentais
  previstos. \\
Regime de opera\c{c}\~ao & Cont\'inuo, sob demanda ou pulsado; fator de
  capacidade estimado. \\
Invent\'ario de produtos de fiss\~ao & Base assumida para o calor de
  decaimento e para os efluentes radioativos liberados. \\
Limita\c{c}\~oes de opera\c{c}\~ao & Restri\c{c}\~oes impostas por
  considera\c{c}\~oes de seguran\c{c}a e sua base nas especifica\c{c}\~oes
  t\'ecnicas. \\
\end{tabelaresumo}

\begin{tabelaapendices}{\ref{sec:1.6}}{Resumo das opera\c{c}\~oes}
Cap\'itulos 4 e 5 -- Reator e Sistemas de Refrigera\c{c}\~ao & Confirmar
  efeito das opera\c{c}\~oes propostas sobre pot\^encia e remo\c{c}\~ao de
  calor. & Pendente \\
Cap\'itulo 14 -- Especifica\c{c}\~oes T\'ecnicas & Confirmar consist\^encia
  entre as limita\c{c}\~oes operacionais citadas e as especifica\c{c}\~oes
  t\'ecnicas propostas. & Pendente \\
\end{tabelaapendices}

\begin{avaliacaointerim}
Aplica-se a mesma amplia\c{c}\~ao de escopo do bloco \ref{sec:1.4}--\ref{sec:1.8}: o resumo
das opera\c{c}\~oes planejadas deve, quando aplic\'avel, cobrir tamb\'em
o regime operacional da instala\c{c}\~ao de produ\c{c}\~ao de
radiois\'otopos (cont\'inuo ou por batelada, conforme o processo de
separa\c{c}\~ao adotado), e n\~ao apenas o do reator. N\~ao h\'a crit\'erio
de aceita\c{c}\~ao adicional nem mudan\c{c}a estrutural motivada pelo
ISG-NUREG-1537 nesta se\c{c}\~ao.
\end{avaliacaointerim}

% =====================================================================
\nusecbreak{1.7}{Se\c{c}\~ao 1.7}
\section{Conformidade com a Gest\~ao de Combust\'ivel Irradiado e
Rejeitos de Alta Atividade}
\label{sec:1.7}

\begin{notaregulatoria}{Se\c{c}\~ao espec\'ifica da legisla\c{c}\~ao dos EUA}
O texto original da NUREG-1537, Se\c{c}\~ao~1.7 \citep{nureg1537p1,nureg1537p2},
trata exclusivamente da conformidade com a \emph{Nuclear Waste Policy Act}
de 1982 dos Estados Unidos, que exige do requerente um contrato com o
\emph{Department of Energy} (DOE) para a destina\c{c}\~ao final de
combust\'ivel irradiado e rejeitos de alta atividade. \textbf{Essa
exig\^encia n\~ao se aplica, tal como redigida, a projetos fora dos
Estados Unidos.} Para um RPAS submetido no Brasil (ou em outra
jurisdi\c{c}\~ao), esta se\c{c}\~ao deve ser reescrita para tratar do
arcabou\c{c}o nacional aplic\'avel \`a gest\~ao de combust\'ivel
irradiado e rejeitos radioativos de alta atividade (no Brasil,
tipicamente envolvendo a CNEN e a pol\'itica nacional de gerenciamento
de rejeitos radioativos, Lei n\textordmasculine~10.308/2001, e a estrutura
da CNEN/ABDAN). Mant\'em-se abaixo a estrutura NUREG apenas como
refer\^encia de formato -- \textbf{o conte\'udo normativo citado n\~ao
deve ser usado literalmente fora do contexto dos EUA.}
\end{notaregulatoria}

[Apenas para requerentes nos EUA:] O requerente discute brevemente como
atende aos requisitos da Se\c{c}\~ao~302(b)(1)(B) da \emph{Nuclear Waste
Policy Act} de 1982 \citep{nwpa1982} para a destina\c{c}\~ao de rejeitos
radioativos de alta atividade e combust\'ivel nuclear irradiado,
incluindo o contrato firmado com o \emph{Department of Energy} (DOE) dos
EUA para a devolu\c{c}\~ao do material. Uma c\'opia da carta de
encaminhamento do contrato entre o requerente e o DOE \'e inclu\'ida em
um ap\^endice deste RPAS (ver Ap\^endice~1.2 -- o documento real firmado
pelo requerente, e n\~ao um exemplo extra\'ido da NUREG-1312; o exemplo da
NUREG-1312 \citep{nureg1312} corresponde ao Ap\^endice~1.1, ver
Se\c{c}\~ao ~\ref{sec:1.1}).

[Para requerentes fora dos Estados Unidos:] Esta se\c{c}\~ao deve, em vez
disso, descrever como o requerente atende ao arcabou\c{c}o nacional
equivalente que rege o gerenciamento, armazenamento e destina\c{c}\~ao
final do combust\'ivel irradiado e dos rejeitos radioativos de alta
atividade [p.ex., no Brasil, a pol\'itica nacional de gerenciamento de
rejeitos radioativos no \^ambito regulat\'orio da CNEN], incluindo
eventuais acordos de armazenamento, licen\c{c}as ou arranjos de
dep\'osito nacional aplic\'aveis.
\begin{tabelaaceitacao}{\ref{sec:1.7}}{Gest\~ao de combust\'ivel irradiado e rejeitos}
1 & (EUA) O requerente deve ter submetido um resumo do contrato com o
     DOE para a destina\c{c}\~ao de rejeitos de alta atividade e
     combust\'ivel irradiado. \\
2 & (EUA) O requerente deve ter indicado onde uma c\'opia da carta do
     contrato pode ser encontrada no RPAS. \\
3 & (Fora dos EUA -- adaptado) O requerente deve identificar o arcabou\c{c}o
     nacional aplic\'avel e demonstrar conformidade com ele. \\
\end{tabelaaceitacao}

\begin{tabelaresumo}{\ref{sec:1.7}}{Gest\~ao de combust\'ivel irradiado e rejeitos}
Arcabou\c{c}o legal aplic\'avel & Norma nacional (ou, nos EUA, a NWPA de
  1982) que rege a destina\c{c}\~ao de combust\'ivel irradiado e rejeitos
  de alta atividade. \\
Arranjo contratual/institucional & Contrato ou acordo com a entidade
  respons\'avel (DOE nos EUA; \'org\~ao nacional equivalente em outros
  pa\'ises). \\
Localiza\c{c}\~ao da evid\^encia documental & Refer\^encia ao ap\^endice
  onde a documenta\c{c}\~ao comprobat\'oria est\'a inclu\'ida. \\
\end{tabelaresumo}

\begin{tabelaapendices}{\ref{sec:1.7}}{Gest\~ao de combust\'ivel irradiado e rejeitos}
Ap\^endice 1.2 (modelo NUREG-1537) -- Carta do DOE & Verificar se este
  ap\^endice deve ser substitu\'ido por documento equivalente da
  jurisdi\c{c}\~ao aplic\'avel. & Pendente -- requer adapta\c{c}\~ao \\
Nuclear Waste Policy Act de 1982, \S302(b)(1)(B) \citep{nwpa1982} &
  Confirmar aplicabilidade; substituir por legisla\c{c}\~ao nacional
  quando fora dos EUA. & Pendente -- requer adapta\c{c}\~ao \\
\end{tabelaapendices}

\pagebreak
\begin{avaliacaointerim}
O ISG-NUREG-1537 n\~ao acrescenta conte\'udo espec\'ifico a esta se\c{c}\~ao
al\'em da amplia\c{c}\~ao gen\'erica de escopo. Vale reforçar, no entanto,
que a ressalva j\'a registrada no corpo desta se\c{c}\~ao -- de que o
texto original da NUREG-1537 trata de uma exig\^encia espec\'ifica da
legisla\c{c}\~ao dos Estados Unidos (a \emph{Nuclear Waste Policy Act} de
1982) sem equival\^encia direta fora daquele pa\'is -- permanece
integralmente v\'alida tamb\'em sob o ISG-NUREG-1537, que n\~ao introduz nenhum
dispositivo alternativo para jurisdi\c{c}\~oes fora dos EUA. A
adapta\c{c}\~ao ao arcabou\c{c}o nacional brasileiro, j\'a descrita no
corpo da se\c{c}\~ao, continua sendo responsabilidade do requerente, n\~ao
uma orienta\c{c}\~ao derivada do ISG-NUREG-1537.
\end{avaliacaointerim}

% =====================================================================
\nusecbreak{1.8}{Se\c{c}\~ao 1.8}
\section{Modifica\c{c}\~oes e Hist\'orico da Instala\c{c}\~ao}
\label{sec:1.8}

Esta se\c{c}\~ao aplica-se principalmente a instala\c{c}\~oes existentes
que solicitam renova\c{c}\~ao de licen\c{c}a; sua aplicabilidade \'e
limitada em uma solicita\c{c}\~ao inicial de permiss\~ao de
constru\c{c}\~ao e licen\c{c}a de opera\c{c}\~ao. [Para uma nova instala\c{c}\~ao:] o requerente apresenta um breve
hist\'orico das atividades pr\'evias relevantes, incluindo eventual
experi\^encia com outros reatores n\~ao-de-pot\^encia ou de pot\^encia.
[Para uma instala\c{c}\~ao em renova\c{c}\~ao ou emenda:] o requerente
apresenta um breve hist\'orico da instala\c{c}\~ao, incluindo as datas de
eventos significativos, a emiss\~ao da permiss\~ao de constru\c{c}\~ao e
da licen\c{c}a de opera\c{c}\~ao, e a criticalidade inicial; e indica se
a instala\c{c}\~ao passou por modifica\c{c}\~oes f\'isicas ou
operacionais significativas ou relacionadas \`a seguran\c{c}a desde o
licenciamento inicial ou desde a \'ultima renova\c{c}\~ao. As
modifica\c{c}\~oes s\~ao discutidas brevemente em ordem cronol\'ogica,
incluindo o n\'umero e a data da emenda de licen\c{c}a aplic\'avel, e as
altera\c{c}\~oes realizadas sob disposi\c{c}\~oes equivalentes \`a
Se\c{c}\~ao~50.59 do 10~CFR~50 \citep{cfr10p50} que afetam as descri\c{c}\~oes
do RPAS s\~ao registradas, juntamente com eventuais altera\c{c}\~oes nas
especifica\c{c}\~oes t\'ecnicas.

\begin{observacao}{Uso desta se\c{c}\~ao para o hist\'orico do RPAS}
Como o RMB \'e uma instala\c{c}\~ao nova, ele n\~ao tem hist\'orico
operacional pr\'evio a relatar aqui (ver caixa de avalia\c{c}\~ao ao
final desta se\c{c}\~ao). Por isso, esta se\c{c}\~ao ser\'a tamb\'em
utilizada para registrar o hist\'orico de elabora\c{c}\~ao e revis\~ao
do pr\'oprio RPAS -- o documento tem passado por diversas altera\c{c}\~oes
ao longo do seu desenvolvimento, como a revis\~ao em curso neste
momento. Esse hist\'orico de vers\~oes ser\'a consolidado e apresentado
aqui \`a medida que as altera\c{c}\~oes forem se acumulando.
\end{observacao}

\begin{tabelaaceitacao}{\ref{sec:1.8}}{Modifica\c{c}\~oes e hist\'orico}
1 & O requerente deve submeter um hist\'orico completo da instala\c{c}\~ao
     (quando aplic\'avel a renova\c{c}\~ao/emenda) ou das atividades
     pr\'evias relevantes (quando instala\c{c}\~ao nova). \\
\end{tabelaaceitacao}

\begin{tabelaresumo}{\ref{sec:1.8}}{Modifica\c{c}\~oes e hist\'orico}
Marcos hist\'oricos & Datas de emiss\~ao de permiss\~oes, licen\c{c}as,
  criticalidade inicial e renova\c{c}\~oes (instala\c{c}\~oes existentes). \\
Modifica\c{c}\~oes significativas & Lista cronol\'ogica de altera\c{c}\~oes
  f\'isicas/operacionais relevantes \`a seguran\c{c}a, com refer\^encia \`a
  emenda de licen\c{c}a correspondente. \\
Altera\c{c}\~oes sob 10~CFR~50.59 (ou equivalente) & Mudan\c{c}as que
  n\~ao exigiram aprova\c{c}\~ao pr\'evia do \'org\~ao regulador, mas que
  afetam as descri\c{c}\~oes do RPAR. \\
\end{tabelaresumo}

\begin{tabelaapendices}{\ref{sec:1.8}}{Modifica\c{c}\~oes e hist\'orico}
Docket regulat\'orio da instala\c{c}\~ao & Confirmar consist\^encia entre
  o hist\'orico apresentado e o registro oficial mantido pelo \'org\~ao
  regulador. & Pendente \\
10~CFR~50.59 (ou disposi\c{c}\~ao nacional equivalente) \citep{cfr10p50} &
  Verificar aplicabilidade do mecanismo de altera\c{c}\~ao sem
  aprova\c{c}\~ao pr\'evia na jurisdi\c{c}\~ao do projeto. & Pendente --
  requer verifica\c{c}\~ao \\
\end{tabelaapendices}

\begin{avaliacaointerim}
Aplica-se a mesma amplia\c{c}\~ao gen\'erica de escopo do bloco \ref{sec:1.4}--\ref{sec:1.8},
sem crit\'erio de aceita\c{c}\~ao adicional espec\'ifico. Como o RMB \'e
uma instala\c{c}\~ao nova (n\~ao uma renova\c{c}\~ao ou emenda de licen\c{c}a
existente), a aplicabilidade desta se\c{c}\~ao j\'a \'e limitada pela
pr\'opria NUREG-1537, e o ISG-NUREG-1537 n\~ao altera essa limita\c{c}\~ao nem
introduz exig\^encia adicional para o caso de instala\c{c}\~ao nova.
\end{avaliacaointerim}

\pagebreak
% =====================================================================


\section*{Ap\^endices citados neste cap\'itulo}
\addcontentsline{toc}{section}{Ap\^endices citados neste cap\'itulo}

\textbf{Ap\^endice~1.1} -- Exemplo de introdu\c{c}\~ao de um relat\'orio
de avalia\c{c}\~ao de seguran\c{c}a (\textit{safety evaluation report}),
reproduzido da NUREG-1312, \emph{``Safety Evaluation Report Related to
the Renewal of the Facility License for the Research Reactor at the Dow
Chemical Company''} \citep{nureg1312}, conforme sugerido pela NUREG-1537,
Parte~2, \S1.1 \citep{nureg1537p2}.

\textbf{Ap\^endice~1.2} -- [Apenas para requerentes nos EUA] C\'opia da
carta do DOE confirmando o arranjo contratual para a destina\c{c}\~ao de
combust\'ivel irradiado e rejeitos de alta atividade sob a \emph{Nuclear
Waste Policy Act} de 1982 \citep{nwpa1982}; deve ser substitu\'ida, fora
dos Estados Unidos, pela documenta\c{c}\~ao nacional equivalente discutida
na Se\c{c}\~ao~\ref{sec:1.7} deste cap\'itulo.

Al\'em desses dois ap\^endices-exemplo, herdados diretamente da estrutura
da NUREG-1537, este RPAS re\'une tr\^es ap\^endices estruturais pr\'oprios,
apresentados ao final do documento: o \textbf{\apref{apA}{A}}
(lista de verifica\c{c}\~ao consolidada dos crit\'erios de aceita\c{c}\~ao
deste cap\'itulo, incluindo os introduzidos pelo ISG-NUREG-1537), o
\textbf{\apref{apB}{B}} (legisla\c{c}\~ao federal do 10~CFR
aplic\'avel, nos moldes do Ap\^endice~A da NUREG-1537) e o
\textbf{\apref{apC}{C}} (discuss\~ao sobre condi\c{c}\~oes
al\'em da base de projeto).
